\documentclass[11pt,a4paper]{article}
\usepackage[margin=25mm]{geometry}
\usepackage[T1]{fontenc}
\usepackage{lmodern,booktabs,tabularx,xcolor,enumitem}
\usepackage[hidelinks]{hyperref}
\usepackage{xurl}
\setlength{\parindent}{0pt}
\setlength{\parskip}{7pt}
\setlength{\emergencystretch}{2em}
% AUTO-GENERATED from actual ten-seed results.
\newcommand{\ExtSimpleBFone}{0.560}
\newcommand{\ExtSimpleBAP}{0.491}
\newcommand{\ExtMultiFone}{0.497}
\newcommand{\ExtMultiReview}{0.365}

\newcommand{\commentbox}[1]{\par\medskip\noindent\colorbox{blue!6}{\parbox{\dimexpr\linewidth-2\fboxsep}{\textbf{Reviewer comment.} #1}}\par}
\begin{document}
\begin{center}
{\LARGE\bfseries Response to Reviewer}\par
\medskip
{\large Annals of Data Science --- Major Revision}\par
\medskip
\textbf{Anomaly Detection Algorithm for Enhanced Fraud Identification in the Insurance Sector}
\end{center}
Dear Editor and Reviewer,

Thank you for the constructive comments. We have applied the revisions
within the original manuscript. Its five main sections, the preprocessing
algorithm, the simple SAD algorithm, the multilayered SAD algorithm and the
original detector/model families are retained. The changes clarify the
contribution, reorganize the literature review and extend the experimental
evaluation. Necessary implementation corrections are disclosed below.
The package includes a clean manuscript, a version marking added and
substantially rewritten passages in blue, a detailed change map, executable
analysis code, and prediction-level outputs. The marked copy is a
reading aid; the change map also identifies removed claims and maps the
original sections and Algorithms 1--3 to their revised counterparts.

\commentbox{The manuscript must more sharply delineate the core innovations
of SAD. The novelty of its ``sequential'' and ``multilayer'' structures is
insufficiently highlighted; a comparative table versus state-of-the-art
counterparts is recommended.}

\textbf{Response.} Sections 1 and 3 now distinguish the mechanisms precisely.
Sequential means that verified records from one batch augment training
before predicting the next batch. Multilayer means that entropy is a gate
applied after anomaly consensus, with entropy recomputed on the candidate
subset. Neither mechanism is presented as a temporal sequence model or a
reinforcement-learning policy.

The contribution statement identifies the review/update workflow and its
empirical tests. The literature comparison table distinguishes supervised,
unsupervised, active/reinforcement and hybrid/graph approaches by learning
signal, strength and boundary. A second comparison table distinguishes
SAD's mechanisms from recent optimized boosting, unsupervised deep learning,
costly-verification RL and claims-graph methods. The two original workflow
images are retained unchanged, and Algorithms 1--3 remain in the methodology
section. An adjacent note clarifies how the executed protocol instantiates
the conceptual diagrams. The submitted article's page layout, title block,
section organization and ggplot2 visual conventions are preserved.
Section 3 supplies the entropy definition, thresholds, training rules,
feature handling and workload accounting. We present the contribution as a
testable integration of existing components, not invention of the individual
detectors or proof of universal superiority.

\commentbox{The literature review needs restructuring. The list-like survey
should be reorganized into supervised, unsupervised,
reinforcement-learning-based, and hybrid architectures, with the strengths,
limitations, and research gaps of each class summarized explicitly.}

\textbf{Response.} Section 2 has been rewritten into the requested four
categories, followed by a synthesis of the research gap and comparator
scope. Each category states its learning signal, practical benefit and
limitation. We distinguish active acquisition from reinforcement learning
and explain that SAD uses fixed review rules without a learned reward policy.

The bibliography now contains 40 cited references, including the method
foundations needed by the restored detector and learner descriptions.
The discussion incorporates recent insurance-specific work on systematic
reviews, imbalance, optimized boosting and claims graphs. Three relevant
\emph{Annals of Data Science} sources support the discussion of imbalanced
feature selection (Kamalov et al., 2023), clustering (Xu and Tian, 2015), and
financial-domain deep learning (Hassani et al., 2020). Their roles are
explicit: the banking review is contextual, not an insurance-fraud numerical
benchmark, and our Gini screen is not the supervised method of Kamalov et al.

\commentbox{The experimental evaluation is incomplete. Recent
fraud-detection-specific baselines are missing, and no ablation study or
validation on an independent second dataset is provided; adding these
experiments would strengthen generalizability.}

\textbf{Response.} Section 4 now reports new executed experiments on two
benchmarks with ten fixed random seeds (1234--1243).

\textbf{Additional empirical comparators.} In addition to the eight component
classifiers, we run balanced random forest and an adapted PSO-XGBoost
comparator motivated by Ding et al. (2025). All classifiers use the same
Simple-B acquired training records. PSO searches only within verified
initial-partition records, with fold-specific SMOTE, and its chosen
hyperparameters remain fixed in later batches. Section 4.4 gives the search
budget and parameter bounds. The comparator is labeled ``adapted'' because
we do not reproduce all data and implementation choices of the publication.
This broadens the evidence while retaining a clear boundary: it is not an
exhaustive reproduction of recent deep or graph-based systems.

\textbf{Ablation and sensitivity.} The Simple-A, Simple-B and Multilayer
comparison studies entropy placement. A new Multilayer-static arm uses the
identical initial reviewed records and fitted models but omits all later
retraining. Its predictions use the same later claims, allowing paired
seed-wise comparisons with updating Multilayer. A second complete
Multilayer run uses $(T_{AD},T_{ML})=(5,6)$ instead of $(3,4)$.
These two joint settings are described as sensitivity analysis, not an
exhaustive grid or an isolated estimate of each threshold's effect.

\textbf{Second benchmark.} The supplied thesis archive contains Insurance
Claims (1,000 records, 38 predictors, 247 recorded frauds), in addition to
Carclaims (15,420 records, 32 predictors, 923 recorded frauds).
The researcher identifies the source of the second file as the
\emph{Insurance Claims} dataset hosted by mexwell on Kaggle; this source
is now cited in Section 4.1 and the bibliography. The analyzed CSV remains
the supplied archival version. The same preprocessing rules, architecture
definitions and SAD hyperparameters are used on both; preprocessing and models are fitted
separately within each benchmark. The differently named
\texttt{fraud\_oracle.csv} is a duplicate representation of Carclaims and is
not counted as an additional dataset.

On Insurance Claims, Simple-B has mean fraud F1 \ExtSimpleBFone{} and AP
\ExtSimpleBAP{}; Multilayer has mean F1 \ExtMultiFone{} and total review
fraction \ExtMultiReview{}. All revised scores, confusion counts and plots
are regenerated from the new runs, rather than copied from legacy exports.
Mean and sample SD across seeds are reported; we do not label them as
independent-sample confidence intervals.

The second file has a different schema and record population, supporting
workflow replication across benchmarks. We cannot establish the original
insurer or real-versus-synthetic provenance of every field from the archive.
We therefore avoid presenting this as prospective validation on a verified
independent insurer portfolio. Section 5 states this remaining limitation,
along with retrospective splitting and previously examined benchmark data.

\textbf{Corrections arising from the reproducibility audit.} The new
implementation excludes the outcome by name before every detector, aligns
row indices across methods, fixes candidate-to-original-row mapping and
entropy gain, and forms acquisition masks before revealing outcomes.
Training uses only verified records, including randomly audited non-alerts.
SMOTE is implemented within training data, in accordance with the original
method description. The revised pseudocode distinguishes accumulated real
verified observations from newly oversampled training data and places
retraining after the complete batch. Fraud is explicitly
the positive class. We distinguish alert rate from total review rate and
report precision, recall, F1, specificity, balanced accuracy, ROC-AUC and AP.
Stored warnings and computation times are retained without selectively
discarding unsuccessful-looking split results.

The executed revised model implementation is Python. The empirical figures
are rendered from those saved results using an executed R/ggplot2 script,
with the same panel, colour, line and point conventions for both datasets.
The original R scripts and outputs are archived for traceability. Each
empirical figure has plot data, generation code, vector exports and a
720-dpi PNG. The two original workflow PNGs and their editable draw.io
sources are retained unchanged. A numerical audit checks
source-label alignment, split and review accounting, and independently
recomputes the reported metrics from distributed prediction files.
Figure labels and axis ranges have been adjusted for readability without
changing numerical values. Section 4.6 now provides the public GitHub
repository containing the data, executable analysis and plotting code,
prediction-level outputs, figure assets and reproduction instructions:
\url{https://github.com/Mehrab-Atighi/Anomaly-Detection-Algorithm-for-Enhanced-Fraud-Identification-in-the-Insurance-Sector}.
Section 3.4 briefly
acknowledges ChatGPT assistance with manuscript editing and revision of the
analysis and plotting scripts; the authors retain responsibility for the
code, results and final text.

Following the editorial instruction, the submission source archive contains
editable LaTeX, bibliography files and PNG figure assets, with no PDF files.
Compiled reading copies and the full computational archive are provided
separately within the complete research package.

We hope the revised presentation and reproducible evidence address the
reviewer's concerns while accurately describing the remaining limits.

Sincerely,\\
Mehrab Atighi, Sayed Jamal Mirkamali, and Amir T. Payandeh Najafabadi
\end{document}
